% Diagrama de clases de la API (PDP) — lámina C: IAM, autenticación, deduplicación y notificaciones
% Compilar: pdflatex 04c_clases_api_iam_auth.tex
\documentclass[tikz,border=6pt]{standalone}
% Estilos comunes de los diagramas (TikZ puro, sin tikz-uml).
% Se carga con % Estilos comunes de los diagramas (TikZ puro, sin tikz-uml).
% Se carga con % Estilos comunes de los diagramas (TikZ puro, sin tikz-uml).
% Se carga con \input{estilos_uml} desde cada diagrama standalone.
\usepackage[T1]{fontenc}
\usepackage[utf8]{inputenc}
\usepackage{lmodern}
\usepackage{underscore}   % permite escribir identificadores con "_" sin escapar
\usepackage{xcolor}
\DeclareUnicodeCharacter{2192}{\ensuremath{\rightarrow}}
\DeclareUnicodeCharacter{2190}{\ensuremath{\leftarrow}}
\DeclareUnicodeCharacter{2194}{\ensuremath{\leftrightarrow}}
\DeclareUnicodeCharacter{2265}{\ensuremath{\geq}}
\DeclareUnicodeCharacter{2264}{\ensuremath{\leq}}
\DeclareUnicodeCharacter{2260}{\ensuremath{\neq}}
\DeclareUnicodeCharacter{00D7}{\ensuremath{\times}}
\DeclareUnicodeCharacter{21D2}{\ensuremath{\Rightarrow}}
\DeclareUnicodeCharacter{2026}{\ldots}
\usetikzlibrary{shapes.multipart,shapes.geometric,shapes.misc,arrows.meta,positioning,fit,backgrounds,calc}

% Paleta (sobria, imprime bien en B/N)
\definecolor{cPep}{HTML}{E3EEF9}     % agente / PEP
\definecolor{cPdp}{HTML}{FCEBD9}     % API / PDP
\definecolor{cDato}{HTML}{E4F2E2}    % datos / persistencia
\definecolor{cExt}{HTML}{EEE6F5}     % servicios externos
\definecolor{cUi}{HTML}{FFF6CC}      % frontend
\definecolor{cOff}{HTML}{EDEDED}     % offline
\definecolor{cBorde}{HTML}{3A3A3A}

\tikzset{
  % --- UML: clase de 3 compartimentos
  clase/.style={rectangle split, rectangle split parts=3, draw=cBorde, fill=white,
    font=\ttfamily\scriptsize, inner sep=2.5pt, text width=#1,
    rectangle split part align={center,left,left}},
  clase/.default=6cm,
  % --- UML: objeto (2 compartimentos, nombre subrayado)
  objeto/.style={rectangle split, rectangle split parts=2, draw=cBorde, fill=white,
    font=\ttfamily\scriptsize, inner sep=2.5pt, text width=#1,
    rectangle split part align={center,left}},
  objeto/.default=5cm,
  % --- relaciones UML
  herencia/.style={-{Triangle[open,length=7pt,width=8pt]}, draw=cBorde, thick},
  realiza/.style={-{Triangle[open,length=7pt,width=8pt]}, draw=cBorde, thick, dashed},
  compone/.style={{Diamond[length=9pt,width=6pt]}-{Stealth[length=5pt]}, draw=cBorde, thick},
  agrega/.style={{Diamond[open,length=9pt,width=6pt]}-{Stealth[length=5pt]}, draw=cBorde, thick},
  usa/.style={-{Stealth[length=5pt]}, draw=cBorde, dashed},
  asocia/.style={-{Stealth[length=5pt]}, draw=cBorde},
  enlace/.style={draw=cBorde, thick},
  rotulo/.style={font=\ttfamily\tiny, fill=white, inner sep=1pt},
  % --- componentes / flujo
  comp/.style={draw=cBorde, rounded corners=2pt, align=center, font=\small, inner sep=4pt, minimum height=9mm},
  paso/.style={draw=cBorde, rounded corners=3pt, align=center, font=\footnotesize, inner sep=3pt, text width=#1},
  paso/.default=3.4cm,
  decision/.style={draw=cBorde, diamond, aspect=2.2, align=center, font=\footnotesize, inner sep=1pt},
  almacen/.style={draw=cBorde, cylinder, shape border rotate=90, aspect=0.18, align=center,
    font=\footnotesize, minimum width=2.6cm, minimum height=1.2cm},
  zona/.style={draw=cBorde, dashed, rounded corners=4pt, inner sep=8pt},
  zonatit/.style={font=\bfseries\small, anchor=north west, inner sep=2pt},
  flujo/.style={-{Stealth[length=6pt]}, draw=cBorde, thick},
  flujofallo/.style={-{Stealth[length=6pt]}, draw=red!60!black, thick, dashed},
  proto/.style={font=\scriptsize\ttfamily, fill=white, inner sep=1pt, align=center},
}

% \umlclase[opciones]{nombre-nodo}{x,y}{estereotipo}{Nombre}{atributos}{operaciones}
% atributos/operaciones: líneas separadas con \\ ; usar \vacio si no hay.
\newcommand{\vacio}{\phantom{x}}
\newcommand{\umlclase}[7][]{%
  \node[clase,#1] (#2) at (#3) {\makebox[\linewidth][c]{%
    \begin{tabular}{@{}c@{}}\if\relax\detokenize{#4}\relax\else\textit{«#4»}\\\fi\textbf{\normalsize #5}\end{tabular}}%
    \nodepart{two}\begin{tabular}{@{}l@{}}#6\end{tabular}%
    \nodepart{three}\begin{tabular}{@{}l@{}}#7\end{tabular}};}

% \umlobjeto[opciones]{nombre-nodo}{x,y}{nombre:Clase}{slots}
\newcommand{\umlobjeto}[5][]{%
  \node[objeto,#1] (#2) at (#3) {\makebox[\linewidth][c]{\underline{\textbf{#4}}}%
    \nodepart{two}\begin{tabular}{@{}l@{}}#5\end{tabular}};}
 desde cada diagrama standalone.
\usepackage[T1]{fontenc}
\usepackage[utf8]{inputenc}
\usepackage{lmodern}
\usepackage{underscore}   % permite escribir identificadores con "_" sin escapar
\usepackage{xcolor}
\DeclareUnicodeCharacter{2192}{\ensuremath{\rightarrow}}
\DeclareUnicodeCharacter{2190}{\ensuremath{\leftarrow}}
\DeclareUnicodeCharacter{2194}{\ensuremath{\leftrightarrow}}
\DeclareUnicodeCharacter{2265}{\ensuremath{\geq}}
\DeclareUnicodeCharacter{2264}{\ensuremath{\leq}}
\DeclareUnicodeCharacter{2260}{\ensuremath{\neq}}
\DeclareUnicodeCharacter{00D7}{\ensuremath{\times}}
\DeclareUnicodeCharacter{21D2}{\ensuremath{\Rightarrow}}
\DeclareUnicodeCharacter{2026}{\ldots}
\usetikzlibrary{shapes.multipart,shapes.geometric,shapes.misc,arrows.meta,positioning,fit,backgrounds,calc}

% Paleta (sobria, imprime bien en B/N)
\definecolor{cPep}{HTML}{E3EEF9}     % agente / PEP
\definecolor{cPdp}{HTML}{FCEBD9}     % API / PDP
\definecolor{cDato}{HTML}{E4F2E2}    % datos / persistencia
\definecolor{cExt}{HTML}{EEE6F5}     % servicios externos
\definecolor{cUi}{HTML}{FFF6CC}      % frontend
\definecolor{cOff}{HTML}{EDEDED}     % offline
\definecolor{cBorde}{HTML}{3A3A3A}

\tikzset{
  % --- UML: clase de 3 compartimentos
  clase/.style={rectangle split, rectangle split parts=3, draw=cBorde, fill=white,
    font=\ttfamily\scriptsize, inner sep=2.5pt, text width=#1,
    rectangle split part align={center,left,left}},
  clase/.default=6cm,
  % --- UML: objeto (2 compartimentos, nombre subrayado)
  objeto/.style={rectangle split, rectangle split parts=2, draw=cBorde, fill=white,
    font=\ttfamily\scriptsize, inner sep=2.5pt, text width=#1,
    rectangle split part align={center,left}},
  objeto/.default=5cm,
  % --- relaciones UML
  herencia/.style={-{Triangle[open,length=7pt,width=8pt]}, draw=cBorde, thick},
  realiza/.style={-{Triangle[open,length=7pt,width=8pt]}, draw=cBorde, thick, dashed},
  compone/.style={{Diamond[length=9pt,width=6pt]}-{Stealth[length=5pt]}, draw=cBorde, thick},
  agrega/.style={{Diamond[open,length=9pt,width=6pt]}-{Stealth[length=5pt]}, draw=cBorde, thick},
  usa/.style={-{Stealth[length=5pt]}, draw=cBorde, dashed},
  asocia/.style={-{Stealth[length=5pt]}, draw=cBorde},
  enlace/.style={draw=cBorde, thick},
  rotulo/.style={font=\ttfamily\tiny, fill=white, inner sep=1pt},
  % --- componentes / flujo
  comp/.style={draw=cBorde, rounded corners=2pt, align=center, font=\small, inner sep=4pt, minimum height=9mm},
  paso/.style={draw=cBorde, rounded corners=3pt, align=center, font=\footnotesize, inner sep=3pt, text width=#1},
  paso/.default=3.4cm,
  decision/.style={draw=cBorde, diamond, aspect=2.2, align=center, font=\footnotesize, inner sep=1pt},
  almacen/.style={draw=cBorde, cylinder, shape border rotate=90, aspect=0.18, align=center,
    font=\footnotesize, minimum width=2.6cm, minimum height=1.2cm},
  zona/.style={draw=cBorde, dashed, rounded corners=4pt, inner sep=8pt},
  zonatit/.style={font=\bfseries\small, anchor=north west, inner sep=2pt},
  flujo/.style={-{Stealth[length=6pt]}, draw=cBorde, thick},
  flujofallo/.style={-{Stealth[length=6pt]}, draw=red!60!black, thick, dashed},
  proto/.style={font=\scriptsize\ttfamily, fill=white, inner sep=1pt, align=center},
}

% \umlclase[opciones]{nombre-nodo}{x,y}{estereotipo}{Nombre}{atributos}{operaciones}
% atributos/operaciones: líneas separadas con \\ ; usar \vacio si no hay.
\newcommand{\vacio}{\phantom{x}}
\newcommand{\umlclase}[7][]{%
  \node[clase,#1] (#2) at (#3) {\makebox[\linewidth][c]{%
    \begin{tabular}{@{}c@{}}\if\relax\detokenize{#4}\relax\else\textit{«#4»}\\\fi\textbf{\normalsize #5}\end{tabular}}%
    \nodepart{two}\begin{tabular}{@{}l@{}}#6\end{tabular}%
    \nodepart{three}\begin{tabular}{@{}l@{}}#7\end{tabular}};}

% \umlobjeto[opciones]{nombre-nodo}{x,y}{nombre:Clase}{slots}
\newcommand{\umlobjeto}[5][]{%
  \node[objeto,#1] (#2) at (#3) {\makebox[\linewidth][c]{\underline{\textbf{#4}}}%
    \nodepart{two}\begin{tabular}{@{}l@{}}#5\end{tabular}};}
 desde cada diagrama standalone.
\usepackage[T1]{fontenc}
\usepackage[utf8]{inputenc}
\usepackage{lmodern}
\usepackage{underscore}   % permite escribir identificadores con "_" sin escapar
\usepackage{xcolor}
\DeclareUnicodeCharacter{2192}{\ensuremath{\rightarrow}}
\DeclareUnicodeCharacter{2190}{\ensuremath{\leftarrow}}
\DeclareUnicodeCharacter{2194}{\ensuremath{\leftrightarrow}}
\DeclareUnicodeCharacter{2265}{\ensuremath{\geq}}
\DeclareUnicodeCharacter{2264}{\ensuremath{\leq}}
\DeclareUnicodeCharacter{2260}{\ensuremath{\neq}}
\DeclareUnicodeCharacter{00D7}{\ensuremath{\times}}
\DeclareUnicodeCharacter{21D2}{\ensuremath{\Rightarrow}}
\DeclareUnicodeCharacter{2026}{\ldots}
\usetikzlibrary{shapes.multipart,shapes.geometric,shapes.misc,arrows.meta,positioning,fit,backgrounds,calc}

% Paleta (sobria, imprime bien en B/N)
\definecolor{cPep}{HTML}{E3EEF9}     % agente / PEP
\definecolor{cPdp}{HTML}{FCEBD9}     % API / PDP
\definecolor{cDato}{HTML}{E4F2E2}    % datos / persistencia
\definecolor{cExt}{HTML}{EEE6F5}     % servicios externos
\definecolor{cUi}{HTML}{FFF6CC}      % frontend
\definecolor{cOff}{HTML}{EDEDED}     % offline
\definecolor{cBorde}{HTML}{3A3A3A}

\tikzset{
  % --- UML: clase de 3 compartimentos
  clase/.style={rectangle split, rectangle split parts=3, draw=cBorde, fill=white,
    font=\ttfamily\scriptsize, inner sep=2.5pt, text width=#1,
    rectangle split part align={center,left,left}},
  clase/.default=6cm,
  % --- UML: objeto (2 compartimentos, nombre subrayado)
  objeto/.style={rectangle split, rectangle split parts=2, draw=cBorde, fill=white,
    font=\ttfamily\scriptsize, inner sep=2.5pt, text width=#1,
    rectangle split part align={center,left}},
  objeto/.default=5cm,
  % --- relaciones UML
  herencia/.style={-{Triangle[open,length=7pt,width=8pt]}, draw=cBorde, thick},
  realiza/.style={-{Triangle[open,length=7pt,width=8pt]}, draw=cBorde, thick, dashed},
  compone/.style={{Diamond[length=9pt,width=6pt]}-{Stealth[length=5pt]}, draw=cBorde, thick},
  agrega/.style={{Diamond[open,length=9pt,width=6pt]}-{Stealth[length=5pt]}, draw=cBorde, thick},
  usa/.style={-{Stealth[length=5pt]}, draw=cBorde, dashed},
  asocia/.style={-{Stealth[length=5pt]}, draw=cBorde},
  enlace/.style={draw=cBorde, thick},
  rotulo/.style={font=\ttfamily\tiny, fill=white, inner sep=1pt},
  % --- componentes / flujo
  comp/.style={draw=cBorde, rounded corners=2pt, align=center, font=\small, inner sep=4pt, minimum height=9mm},
  paso/.style={draw=cBorde, rounded corners=3pt, align=center, font=\footnotesize, inner sep=3pt, text width=#1},
  paso/.default=3.4cm,
  decision/.style={draw=cBorde, diamond, aspect=2.2, align=center, font=\footnotesize, inner sep=1pt},
  almacen/.style={draw=cBorde, cylinder, shape border rotate=90, aspect=0.18, align=center,
    font=\footnotesize, minimum width=2.6cm, minimum height=1.2cm},
  zona/.style={draw=cBorde, dashed, rounded corners=4pt, inner sep=8pt},
  zonatit/.style={font=\bfseries\small, anchor=north west, inner sep=2pt},
  flujo/.style={-{Stealth[length=6pt]}, draw=cBorde, thick},
  flujofallo/.style={-{Stealth[length=6pt]}, draw=red!60!black, thick, dashed},
  proto/.style={font=\scriptsize\ttfamily, fill=white, inner sep=1pt, align=center},
}

% \umlclase[opciones]{nombre-nodo}{x,y}{estereotipo}{Nombre}{atributos}{operaciones}
% atributos/operaciones: líneas separadas con \\ ; usar \vacio si no hay.
\newcommand{\vacio}{\phantom{x}}
\newcommand{\umlclase}[7][]{%
  \node[clase,#1] (#2) at (#3) {\makebox[\linewidth][c]{%
    \begin{tabular}{@{}c@{}}\if\relax\detokenize{#4}\relax\else\textit{«#4»}\\\fi\textbf{\normalsize #5}\end{tabular}}%
    \nodepart{two}\begin{tabular}{@{}l@{}}#6\end{tabular}%
    \nodepart{three}\begin{tabular}{@{}l@{}}#7\end{tabular}};}

% \umlobjeto[opciones]{nombre-nodo}{x,y}{nombre:Clase}{slots}
\newcommand{\umlobjeto}[5][]{%
  \node[objeto,#1] (#2) at (#3) {\makebox[\linewidth][c]{\underline{\textbf{#4}}}%
    \nodepart{two}\begin{tabular}{@{}l@{}}#5\end{tabular}};}

\begin{document}
\begin{tikzpicture}

% ===================== FILA 1
\umlclase[clase=10.4cm, fill=cPdp, anchor=north]{iamsvc}{-15,0}{module}{iam_service \textmd{(cache + fail-safe)}}{%
 + IAM_CACHE_TTL_SECONDS = 1800\\
 + IAM_FAIL_CACHE_TTL_SECONDS = env IAM_FAIL_CACHE_TTL | 60\\
 + NIVELES_REVOCAN_SESION = frozenset(\{"critico"\})}{%
 - \_cache_key(uid: str): str\\
 - \_get_cached(uid: str): tuple[str | None, str | None] | None\\
 - \_store_cached(uid: str, role: str | None, department: str | None,\\
 \quad ttl: float = IAM_CACHE_TTL_SECONDS): None\\
 + resolver_rol_departamento(uid: str): tuple[str | None, str | None]\\
 - \_sync_fail_key(uid: str): str\\
 + sincronizar_riesgo_iam(uid: str, nivel: str | None): bool}

\umlclase[clase=10cm, fill=white, anchor=north]{iamp}{0,0}{Protocol}{IAMProvider}{\vacio}{%
 + resolve(uid: str): tuple[str | None, str | None]\\
 + aplicar_riesgo(uid: str, nivel: str, revocar_sesiones: bool): bool}

\umlclase[clase=10.4cm, fill=cPdp, anchor=north]{dedup}{17.6,-38.6}{module}{dedup_service}{%
 + DEDUP_TTL_SECONDS = 21600 \hfill (6 h)\\
 + CAMPOS_HASH = ("uid", "time", "type", "params",\\
 \quad "location", "role", "isLocalIP")}{%
 - \_dedup_key(hash_hex: str): str\\
 + compute_event_hash(event: dict): str \hfill SHA-256\\
 + get_cached_response(hash_hex: str): dict | None\\
 + store_response(hash_hex: str, response: dict): None}

% ===================== FILA 2: implementaciones
\umlclase[clase=8.6cm, fill=cPdp, anchor=north]{mock}{-17.2,-8}{}{MockIAMProvider}{\vacio}{%
 + resolve(uid: str): tuple[str | None, str | None]\\
 \hfill GET \{IAM_MOCK_URL\}/Users/\{uid\}\\
 + aplicar_riesgo(uid: str, nivel: str,\\
 \quad revocar_sesiones: bool): bool \hfill (siempre False)}

\umlclase[clase=10.2cm, fill=cPdp, anchor=north]{kc}{-6.4,-8}{}{KeycloakIAMProvider}{%
 + cfg: \_KcCfg}{%
 + \_\_init\_\_(cfg: \_KcCfg | None = None)\\
 - \_buscar_usuario(uid: str, token: str): dict | None\\
 \hfill GET /admin/realms/\{realm\}/users?username=\&exact=true\\
 + resolve(uid: str): tuple[str | None, str | None]\\
 \hfill attributes.role / attributes.department\\
 + aplicar_riesgo(uid: str, nivel: str, revocar_sesiones: bool): bool\\
 \hfill PUT /users/\{id\} (ueba_risk_level); POST /users/\{id\}/logout}

\umlclase[clase=10.2cm, fill=cPdp, anchor=north]{entra}{4.6,-8}{}{EntraIAMProvider}{%
 + cfg: \_EntraCfg}{%
 + \_\_init\_\_(cfg: \_EntraCfg)\\
 - \_usuario(uid: str, token: str): dict | None\\
 \hfill GET graph /users/\{uid\}?\$select=id,jobTitle,department\\
 + resolve(uid: str): tuple[str | None, str | None]\\
 \hfill (jobTitle, department)\\
 + aplicar_riesgo(uid: str, nivel: str, revocar_sesiones: bool): bool\\
 \hfill PATCH /users/\{id\} extensionAttribute1;\\
 \hfill POST /users/\{id\}/revokeSignInSessions}

\umlclase[clase=8.6cm, fill=cPdp, anchor=north]{chain}{15.4,-8}{}{ChainIAMProvider \textmd{(Composite)}}{%
 + providers: list[IAMProvider]}{%
 + \_\_init\_\_(providers: list[IAMProvider])\\
 + resolve(uid: str): tuple[str | None, str | None]\\
 \hfill primer resultado ≠ (None, None) gana\\
 + aplicar_riesgo(uid: str, nivel: str,\\
 \quad revocar_sesiones: bool): bool}

% ===================== FILA 3
\umlclase[clase=5.4cm, fill=white, anchor=north]{kccfg}{-6.4,-16.2}{NamedTuple}{\_KcCfg}{%
 + base_url: str\\
 + realm: str\\
 + client_id: str\\
 + client_secret: str}{\vacio}

\umlclase[clase=5.4cm, fill=white, anchor=north]{entracfg}{4.6,-16.2}{NamedTuple}{\_EntraCfg}{%
 + tenant: str\\
 + client_id: str\\
 + client_secret: str}{\vacio}

\umlclase[clase=11.4cm, fill=cPdp, anchor=north]{iamprov}{17.6,-19}{module}{iam_providers}{%
 + IAM_MOCK_URL = env | "http://127.0.0.1:8002"\\
 + IAM_TIMEOUT_SEC = 1.0\\
 + KEYCLOAK_BASE_URL = env | "http://127.0.0.1:8080"\\
 + KEYCLOAK_REALM, KEYCLOAK_CLIENT_ID, KEYCLOAK_CLIENT_SECRET: env\\
 + KEYCLOAK_TIMEOUT_SEC = 1.0\\
 + KEYCLOAK_TOKEN_REFRESH_MARGIN_SEC = 30\\
 + ATRIBUTO_RIESGO = "ueba_risk_level"\\
 + GRAPH_BASE_URL = "https://graph.microsoft.com/v1.0"\\
 + ENTRA_LOGIN_URL = "https://login.microsoftonline.com"\\
 + ENTRA_TIMEOUT_SEC = 1.0\\
 + ENTRA_ATRIBUTO_RIESGO = "extensionAttribute1"\\
 + IAM_PROVIDER = env | "none"\\
 + PROVEEDORES: dict[str, Callable[[dict], IAMProvider]]\\
 \quad \{"mock", "keycloak", "entra"\}\\
 - \_iam_provider: IAMProvider | None \hfill «singleton»\\
 - \_config_redis: list[dict] | None}{%
 - \_cfg_default(): \_KcCfg\\
 - \_keycloak_token_cache_key(cfg: \_KcCfg | None): str\\
 - \_get_cached_token(cfg: \_KcCfg | None): str | None\\
 - \_store_cached_token(access_token: str, expires_in: float,\\
 \quad cfg: \_KcCfg | None): None\\
 - \_fetch_service_token(cfg: \_KcCfg | None): tuple[str, float]\\
 \hfill POST /realms/\{realm\}/protocol/openid-connect/token\\
 - \_get_service_token(cfg: \_KcCfg | None): str\\
 - \_entra_token_cache_key(cfg: \_EntraCfg): str\\
 - \_fetch_entra_token(cfg: \_EntraCfg): tuple[str, float]\\
 \hfill POST /\{tenant\}/oauth2/v2.0/token\\
 - \_get_entra_token(cfg: \_EntraCfg): str\\
 - \_build_provider(spec: str): IAMProvider \hfill mock | keycloak[:N] | entra:N\\
 + proveedores_entorno(): list[dict]\\
 - \_config_redis_items(): list[dict]\\
 + reset_provider(): None\\
 + iam_habilitado(): bool\\
 - \_from_item(item: dict): IAMProvider\\
 + get_iam_provider(): IAMProvider\\
 + probar_proveedores(): list[dict] \hfill \{proveedor, ok, ms, error\}}

\umlclase[clase=11.4cm, fill=cPdp, anchor=north]{auth}{-15.6,-19}{module}{auth}{%
 + SESION_TTL_SECONDS = 28800 \hfill (8 h)\\
 + PASSWORDS_PLACEHOLDER = ("TU_PASSWORD_ACA",\\
 \quad "PLACEHOLDER_CARGAR_MANUALMENTE")\\
 - \_NO_APLICA = object()\\
 env: SCORING_API_KEY, SCORING_ADMIN_KEY, SCORING_STRICT_AUTH,\\
 \quad ADMIN_PASSWORD, OPERADOR_PASSWORD}{%
 - \_eq(a: str, b: str): bool \hfill hmac.compare_digest\\
 - \_configured(name: str): str | None\\
 + strict_auth_enabled(): bool\\
 + keys_configuradas(): list[str]\\
 + admin_key_configurada(): bool\\
 + agent_key_configurada(): bool \hfill SCAN agent_key:*\\
 - \_agent_key_lookup(raw_key: str): dict | None\\
 - \_session_key(token: str): str\\
 + crear_sesion(role: str): str \hfill token_urlsafe(32)\\
 - \_rol_de_sesion(token: str): str | None\\
 + cerrar_sesion(token: str): None\\
 + require_api_key(x_api_key = Header(None)): None \hfill «Depends» 401/503\\
 + require_agent_key(x_api_key = Header(None)): dict | None \hfill 401/503\\
 + autorizar_agente(entrada: dict | None, uid: str,\\
 \quad agent_id = \_NO_APLICA): None \hfill 403\\
 + require_admin_key(x_api_key = Header(None)): None \hfill 403/503\\
 - \_password_configurada(var: str, valor_ejemplo: str): str | None\\
 + authenticate(username: str, password: str): str | None \hfill admin | operador}

% ===================== FILA 4: notificaciones
\umlclase[clase=15cm, fill=cPdp, anchor=north]{notif}{1,-28.4}{module}{notifications \textmd{(acción notificar_admin)}}{%
 + SMTP_TIMEOUT_SEC = 3\\
 + NOTIF_COOLDOWN_SEC = env | 600\\
 + GEMINI_API_URL = "https://generativelanguage.googleapis.com/v1beta/models/\{modelo\}:generateContent"\\
 + GEMINI_MODELOS = ["gemini-3.1-flash-lite", "gemini-2.5-flash"]\\
 + GEMINI_TIMEOUT_SEC = 6.0, GEMINI_MAX_REINTENTOS_429 = 1\\
 + MAX_CHARS_CAMPO = 80, MAX_EVENTOS_PROMPT = 10, MAX_CHARS_SALIDA = 1500, MIN_CHARS_SALIDA = 20\\
 + SALIDA_PROHIBIDA, ENCABEZADO_SALIDA: re.Pattern, IDENTIFICADOR_SNAKE: re.Pattern\\
 + TIPO_EVENTO_DESC, FEATURE_DESCRIPCIONES, DIRECCION_LEGIBLE: dict\\
 + GUIA_QUE_HACER_AHORA: list[str], GEMINI_REGLAS: str}{%
 - \_sanear(valor, max_len: int = MAX_CHARS_CAMPO): str\\
 - \_formatear_guia_que_hacer_ahora(): str\\
 - \_hora_legible(timestamp_iso: str | None, con_segundos: bool = True): str | None\\
 - \_formatear_evento_reciente(evento: dict): str\\
 - \_formatear_eventos_recientes_desde(eventos: list[dict]): str\\
 - \_humanizar_feature(nombre_feature: str): str\\
 - \_formatear_explicacion(explicacion: list[dict] | None): str\\
 - \_fmt_umbral(valor: float): str\\
 - \_rango_nivel(nivel: str | None): str | None\\
 - \_armar_datos_gemini(uid: str, nivel: str | None, score: float | None, eventos: list[dict],\\
 \quad explicacion: list[dict] | None): str\\
 - \_validar_salida(texto: str, finish_reason: str | None): str | None\\
 - \_generar_narrativa_gemini(uid, nivel, score, eventos, explicacion): str | None \hfill HTTPS REST\\
 + notificar_admin_si_corresponde(uid: str, nivel: str | None, score: float | None,\\
 \quad explicacion: list[dict] | None = None, timestamp: str | None = None,\\
 \quad accion_id: str | None = None): None \hfill smtplib STARTTLS :587}

\node[draw=cBorde, fill=cUi, font=\ttfamily\scriptsize, align=left, text width=7cm, anchor=north]
 (notadep) at (-15.6,-38.6) {\textbf{Dependencias en otras láminas}\\
 iam_service, iam_providers, auth, dedup_service\\\quad → redis_state (lámina A)\\
 iam_providers → config_store.get_iam_config (lámina D)\\
 notifications → artifacts (B), config_store (D),\\\quad redis_state (A)\\
 hybrid_service (B) invoca iam_service y notifications;\\
 routes (D) invoca auth y dedup_service};

% ===================== RELACIONES
% realización del Protocol (bus)
\coordinate (busP) at (0,-6.6);
\draw[realiza] (0,-6.6) -- (iamp.south);
\draw[draw=cBorde, thick, dashed] (mock.north |- busP) -- (chain.north |- busP);
\foreach \n in {mock, kc, entra, chain}
  \draw[draw=cBorde, thick, dashed] (\n.north) -- (\n.north |- busP);
\node[rotulo, right] at ([xshift=2mm]0,-6.0) {«realiza» (estructural)};

% Composite: Chain agrega 1..* IAMProvider
\draw[agrega] (chain.east) -- ++(0.8,0) |- node[rotulo,pos=0.75,above]{providers 1..*} ([yshift=-0.4cm]iamp.east);

% configuración de cada provider
\draw[compone] (kc.south) -- node[rotulo,right]{cfg 1} (kccfg.north);
\draw[compone] (entra.south) -- node[rotulo,right]{cfg 1} (entracfg.north);

% servicios
\draw[usa] (iamsvc.east |- iamp.west) -- node[rotulo,above]{«usa» get_iam_provider().resolve / aplicar_riesgo} (iamp.west);
\draw[usa] (iamprov.north) -- node[rotulo,right]{«crea» (\_build_provider, \_from_item)} (iamprov.north |- chain.south);

\end{tikzpicture}
\end{document}
